\documentstyle{article}
\begin{document}

\centerline{\bf Using Lie Groups to Solve Differential Equations}

\vskip 0.3cm

\centerline{Stephen V. Coggeshall}
\centerline{XCM, MS F645}
\centerline{Los Alamos National Laboratory}
\centerline{Los Alamos, NM 87545, USA}
\centerline{Email: svc@demos.lanl.gov}

\vskip 0.3cm

Continuous groups, or Lie groups, were developed by Sophus Lie in the latter
part of the 19th Century specifically for the reduction and solution of
differential equations. After his development their use with differential
equations was virtually forgotten, while they were picked up and used in a
number of other applications. Around the middle of this century a resurgence
in the use with differential equations occurred, leading to an
ever-increasing number of applications, monographs and papers.

The property of invariance under a nontrivial Lie group allows (1) first
order ODE's to be reduced to quadrature, (2) Nth order ODE's reduced to
N-1th order, (3) PDE's with 2 independent variables transformed to ODE's,
and (4) PDE's with N independent variables transformed to PDE's with N-1
independent variables. Under certain conditions the techniques can be
applied multiple times to achieve multiple reductions. The methods are valid
for both linear and nonlinear differential equations, and it is found that
many of the particular, ad hoc techniques of the solution of differential
equations have their basis in this theory.

The general concepts of the theory will be discussed, with emphasis on and
examples from PDE's. The conditions for multiple reductions require an
investigation of the algebraic properties of the groups and the process is
facilitated by the development of optimal systems, which will be discussed.
Examples from 3-D hydrodynamics will be shown.

\end{document}
